% Propietario conservado: /Users/ruben/Documents/ChatGPT/jueces y controles/REVISION_ARGUMENTAL_APERTURAS_CIERRES_20260915/II/manuscript_es/sections/06b_significado_incidencia.tex
\subsection{Dépendance structurelle de la fonctionnelle et normalisation aréale}

La construction précédente situe le paramètre de Barbero--Immirzi à la
rencontre de trois déterminations : incidence marquée, orientation
dodécaphasique et évaluation angulaire. L'incidence apporte les espaces
d'événements ; l'orientation fixe la chaîne sur ces espaces ;
l'évaluation transporte la chaîne vers une quantité adimensionnelle. Cette
décomposition permet de conserver l'origine de chaque coefficient lors de
l'utilisation ultérieure de la fonctionnelle dans la géométrie d'aire.

L'inclusion \(\mathcal F_6\hookrightarrow\mathcal F_8\) maintient
la paire marquée de l'hexade. Son complément a pour cardinal
\(2\binom62=30\). Cet incrément réapparaît comme déterminant de
l'application qui reconstruit les occupations du bord à partir de l'aire et
de l'observable pondéré. L'égalité de ces deux entiers possède ici
une explication concrète : les poids sectoriels diffèrent précisément
du nombre d'événements ajoutés par l'inclusion marquée.

\begin{proposition}[Incidence et récupérabilité des occupations]
Soient \(n_6,n_8\) les occupations agrégées de deux secteurs et soient
\(r_6=|\mathcal F_6|\), \(r_8=|\mathcal F_8|\). L'application
\[
 (n_6,n_8)\longmapsto(n_6+n_8,r_6n_6+r_8n_8)
\]
est inversible sur son image lorsque l'inclusion ajoute au moins un
événement. Pour l'inclusion construite, son déterminant vaut 30.
\end{proposition}
\begin{proof}
La matrice de l'application est
\(\left(\begin{smallmatrix}1&1\\r_6&r_8\end{smallmatrix}\right)\).
Son déterminant est \(r_8-r_6=|\mathcal F_8\setminus\mathcal F_6|\).
L'inclusion considérée ajoute \(2\binom62=30\) événements. L'inversion
linéaire restitue les deux occupations et conserve les contraintes entières
propres à son image.
\end{proof}

Le résultat restitue des occupations agrégées. L'information concernant
chaque liaison et l'ordre des événements demeure dans l'état complet
et dans ses lecteurs supplémentaires. Cette hiérarchie donne une formulation
précise de la différence entre valeur d'aire et généalogie de frontière :
l'aire conserve une somme ; le couple d'observables conserve la
décomposition de cette somme par secteurs ; le registre conserve en outre
sa distribution et son histoire.

Le paramètre \(\gamma_{\HMT}\) transporte cette incidence à l'échelle
aréale. Son interprétation structurelle est celle d'un coefficient orienté
de la fonctionnelle de retour hexadique--octadique. La réalisation gravitationnelle
le reconnaît ensuite dans l'action de Holst. On distingue ainsi sa
définition intrinsèque, sa valeur et sa fonction dans l'équation physique.
